% !TEX root = ../main.tex

\chapter{Discussion}
\label{cha:discussion}

The results in Chapter~\ref{cha:results} paint an ambivalent picture that deserves careful discussion.
On one hand, \sculpt is technically successful: the language is Turing-complete, its interpreter works, the physical pieces withstand use, and most participants identify and manipulate the components without difficulty.
On the other hand, writing programs in \sculpt remains hard.
Roughly half of the participants complete the computational tasks, and expert programmers do not outperform novice teenagers on that metric.

This tension is not a failure of the language but one of its constitutive features.
\sculpt does not seek to eliminate the intrinsic difficulty of programming, a problem shared by every language.
The intention is to change the \emph{nature} of that difficulty and, in doing so, to broaden who finds it interesting.
What the results show is that the difficulty is redistributed: the barrier to first interaction with the language falls sharply, children manipulate pieces and artists enjoy them.
What stays almost untouched is the difficulty of reasoning algorithmically about stacks, conditionals and loops.
The discrepancy between the low technical perception of participants (Likert means 2.2--3.4) and the high educator perception (4.0--4.3) reinforces this reading.
\sculpt works better as a didactic vehicle mediated by a facilitator than as an autonomous programming tool.

\section{Physicality and Aesthetics as Objects of Study}
\label{sec:discussion:aesthetics}

Physicality and aesthetics deserve to be treated as objects of study in their own right, not merely as means to an end.
The sessions with artists are the clearest case.
Participants who reported the greatest aesthetic satisfaction were also those who most pulled the pieces away from their prescribed computational use, building non-functional sculptures and exploring assemblages not anticipated in the specification.
In their hands the language stopped being a tool and became a material.
This observation connects directly to open questions in the \emph{programming experience} community: what does it feel like to program in a medium that can be held, what creative purposes emerge when code leaves the screen.

\section{Positioning}
\label{sec:discussion:positioning}

\sculpt does not aspire to replace textual languages.
Its cumbersome nature for large software projects is instantly apparent.
The size of a program is bounded by the number of pieces that can physically be assembled, the printing cost grows linearly with that number, the assembly takes time, and, being only physical, the assembled result does not run on a computer.
\sculpt functions, instead, as an entry point for audiences historically distant from programming and as an expressive medium for artists who want to incorporate computation into their practice without going through the keyboard.

\section{Limitations}
\label{sec:discussion:limitations}

Beyond the size constraints already mentioned, three further limitations are worth naming.
The first and largest is rooted in the physical nature of the exercise.
The validation sessions require facilitators who guide the participants, resolve doubts, help assemble sculptures, supervise the use of the pieces and, in the case of minors, make sure that teachers accompany the students.
This makes the scalability of the exercise limited, and its application in educational contexts dependent on the availability of trained facilitators.

The second limitation goes hand in hand with the first.
The use of 3D-printed pieces with magnetic connectors is a practical solution for prototyping and manufacturing the language at an accessible cost.
The material is fragile and wears down with use, however, reducing the inventory available for each session.
The time cost of printing enough pieces to run the workshops is also high given the limited availability of 3D printers.
This restricts the number of participants who can be served in a session and is visible in the results: the number of participants is small and the number of computational tasks assigned is limited.
The limitation can be overcome with more resistant materials, mass printing processes, and design optimisation.

Third, the validation samples, although diverse, are not large.
Educational performance over the longer term, in stable cohorts, remains an open question.

\endinput
